% ── 다리 초안: allart2018 → 우리 ΔS^0_j 검증 대조(§2.9 procedurebox 4 / tab:ds) ──
% 삽입 위치: ch2_sec09_method.tex, procedurebox(현행 10–29행) 직후,
%   "\subsection{정직한 한계}"(현행 31행) 앞.
% 앵커 문장: "...중심 기울기에서 $\Delta S^0_j=F\,\dd U_j/\dd T$ 를 얻고, Allart 등의 전극 분리
%   $\Delta S(x)$ \cite{allart2018} 와 부호$\cdot$규모를 대조해 검증한다(표~\ref{tab:ds})."(현행 21–22행)
% 컨테이너: srcbox. 우리 기호로 서술. 데이터/량 대응 다리(유도 아님).
% ★확인 필요 없음(논문 식 재현 안 함) — 방법 요지는 논문 제목/주제로 검증. 단 tab:ds 의
%   구체 수치 대조는 원문 값 확인 전 "부호·규모 수준"으로만 한정(수치 재현 금지).

\begin{srcbox}
\textbf{다리 --- $\Delta S^0_j$ 를 무엇으로 검증하나.} procedurebox 4 의 회귀량 $\Delta S^0_j=F\,\dd
U_j/\dd T$ 는 Allart 등\cite{allart2018} 의 흑연 전극 분리 $\Delta S(x)$ 와 대조해 검증한다:
\begin{center}\footnotesize
\begin{tabular}{@{}ll@{}}
\hline
본문 & Allart\cite{allart2018} 대응 \\
\hline
전이별 중심 표준값 $\Delta S^0_j=F\,\dd U_j/\dd T$ & 흑연 전극 몫 $\Delta S(x)$(조성 분해 실측) \\
회귀 $U_j(T)$ 중심 기울기 & 평형 곡선 $U(x)$ 의 온도 의존 \\
표~\ref{tab:ds} 부호$\cdot$규모 & $\Delta S(x)$ 실측 부호$\cdot$규모 \\
\hline
\end{tabular}
\end{center}
\emph{방법 요지} --- Allart 등은 흑연 하프셀의 평형 곡선과 온도계단 전위차법 엔트로피 변화 곡선을 전위차
분석으로 함께 얻어 흑연 전극 몫의 $\Delta S(x)$ 를 조성 분해로 준다. \emph{가정 차} --- Allart 의 $\Delta S(x)$
는 조성 연속의 전극 몫 실측이나 표~\ref{tab:ds} 의 $\Delta S^0_j$ 는 전이별 봉우리 \emph{중심} 표준값이고
봉우리 내부 $x$-의존은 식~\eqref{eq:single_config} 의 config 항이 자동 생성하므로, 대조는 중심값
부호$\cdot$규모 수준이지 봉우리 내부 프로파일의 점대점 일치가 아니다.
\end{srcbox}
